\section{Address Space}

\begin{itemize}
  \item VIV-32 uses a flat 32-bit physical address space.
  \item Addresses are 32-bit unsigned values.
  \item Memory is byte-addressed.
  \item Each byte is 8 bits.
  \item Virtual memory is not supported in VIV-32 v0.1.
  \item Segmentation is not supported in VIV-32 v0.1.
  \item Paging is not supported in VIV-32 v0.1.
\end{itemize}

\section{Endianness}

\begin{itemize}
  \item VIV-32 is big-endian.
  \item Multi-byte integer values are stored with the most significant byte at the lowest memory address.
  \item A 32-bit word value \texttt{0x12345678} stored at address \texttt{A} is represented as:
\end{itemize}

\begin{center}
\begin{tabular}{ll}
\hline
Address & Byte value \\
\hline
\texttt{A + 0} & \texttt{0x12} \\
\texttt{A + 1} & \texttt{0x34} \\
\texttt{A + 2} & \texttt{0x56} \\
\texttt{A + 3} & \texttt{0x78} \\
\hline
\end{tabular}
\end{center}

\section{Access Sizes}

\begin{itemize}
  \item VIV-32 supports 8-bit byte loads and stores.
  \item VIV-32 supports 16-bit halfword loads and stores.
  \item VIV-32 supports 32-bit word loads and stores.
  \item 64-bit ABI values are represented using two 32-bit words or eight bytes of memory.
  \item Native 64-bit memory operations are not required in VIV-32 v0.1.
\end{itemize}

\section{Alignment}

\begin{itemize}
  \item Byte accesses may occur at any address.
  \item Halfword accesses must be aligned to a 2-byte boundary.
  \item Word accesses must be aligned to a 4-byte boundary.
  \item 64-bit ABI values stored in memory should be aligned to an 8-byte boundary.
  \item Misaligned halfword or word accesses raise a misaligned access exception if exception handling is enabled.
  \item If exception handling is not enabled, misaligned access halts the machine.
\end{itemize}

\section{Instruction Alignment}

\begin{itemize}
  \item VIV-32 instructions are 32 bits wide.
  \item Instruction fetch addresses must be 4-byte aligned.
  \item Branch and jump targets must be 4-byte aligned.
  \item Attempting to fetch an instruction from a misaligned address raises a misaligned access exception or halts the machine if exception handling is unavailable.
\end{itemize}

\section{Memory-Mapped I/O}

\begin{itemize}
  \item VIV-32 uses memory-mapped I/O.
  \item Device registers occupy reserved physical address ranges.
  \item Exact device addresses are defined by the platform specification, not by the base ISA.
  \item The architecture reserves the upper physical address range for platform-defined MMIO.
  \item The default platform profile may reserve \texttt{0xFFFF0000--0xFFFFFFFF} for MMIO.
\end{itemize}

\section{Reserved Low Memory}

\begin{itemize}
  \item The platform may reserve low physical memory for reset, vectors, boot code, null-pointer detection, or diagnostic use.
  \item The default platform profile may reserve the first 4 KiB of physical memory.
  \item The exact low-memory layout is defined by the platform specification.
\end{itemize}

\section{Caching}

\begin{itemize}
  \item VIV-32 v0.1 defines no cache architecture.
  \item All memory accesses are treated as direct accesses to memory or memory-mapped devices.
  \item Cache behaviour, cache coherency, and memory barriers are reserved for future versions.
\end{itemize}

\section{Atomicity}

\begin{itemize}
  \item VIV-32 v0.1 does not define atomic read-modify-write memory instructions.
  \item Naturally aligned 8-bit, 16-bit, and 32-bit loads and stores are individually completed as single architectural memory operations.
  \item Multiprocessor atomicity, shared-memory locking, and inter-CPU memory ordering are deferred to a future extension.
\end{itemize}

\section{Memory Ordering}

\begin{itemize}
  \item VIV-32 v0.1 assumes in-order execution.
  \item A single CPU observes its own memory operations in program order.
  \item No multiprocessor memory consistency model is defined in VIV-32 v0.1.
  \item Future versions may define memory barriers, atomic operations, and inter-processor ordering rules.
\end{itemize}
